\appendixitem{连通紧群的结构}

设 G 是交换紧群。回忆（Spectral Theory, Chap. II, §1, no. 9, Prop. 11），此时 G 同构于离散交换群 \(\hat{G}\) 的对偶拓扑群。群 G 连通，当且仅当 \(\hat{G}\) 无挠（Spectral Theory, Chap. II, §2, no. 2, Cor. 1 of Prop. 4）。

下列性质等价（Spectral Theory, Chap. II, §2, no. 2, Cor. 2 of Prop. 4 and §1, no. 9, Cor. 2 of Prop. 11）：

\begin{enumerate}
\def\labelenumi{(\roman{enumi})}
\item
  G 是全不连续的；
\item
  \(\hat{\mathrm{G}}\) 是挠群；
\item
  拓扑群 G 同构于一个有限（交换）群射影系统的极限，其中每个群都带有离散拓扑。
\end{enumerate}

下述命题推广了 §1, no. 4 的命题 4 的推论 1。

命题 2. 设 G 是连通紧群。

\begin{enumerate}
\def\labelenumi{\alph{enumi})}
\item
  \(\mathrm{C}(\mathrm{G})_{0}\) 是交换连通紧群；\(\mathrm{D}(\mathrm{G})\) 是连通紧群，且等于其换位子群。
\item
  连续同态 \((x,y)\mapsto xy\) 从 \(\mathrm{C}(\mathrm{G})_{0}\times\mathrm{D}(\mathrm{G})\) 到 G 是满射，其核是 \(\mathrm{C}(\mathrm{G})_{0}\times\mathrm{D}(\mathrm{G})\) 的中心子群，并且是紧的、全不连续的。
\item
  存在一族概单紧李群 \((\mathrm{S}_{\lambda})_{\lambda\in\mathrm{L}}\) 以及满射连续同态 \(\prod_{\lambda\in\mathrm{L}}\mathrm{S}_{\lambda}\to\mathrm{D}(\mathrm{G})\)，其核是全不连续的紧中心子群。
\end{enumerate}

设 \((\mathrm{G}_{\alpha}, f_{\alpha\beta})\) 是紧李群的射影系统，相对于一个滤集 I，使得 G 同构于 \(\varprojlim G_{\alpha}\)，且典则映射 \(f_{\alpha}: G \to G_{\alpha}\) 是满射（命题 1 的推论 2）。对于 \(\alpha \in I\)，令 \(\pi_{\alpha}: \tilde{\mathrm{D}}(\mathrm{G}_{\alpha}) \to \mathrm{D}(\mathrm{G}_{\alpha})\) 为群 \(\mathrm{D}(\mathrm{G}_{\alpha})\) 的泛覆盖。映射 \(f_{\alpha\beta}\) 诱导出态射 \(\tilde{f}_{\alpha\beta}: \tilde{\mathrm{D}}(\mathrm{G}_{\beta}) \to \tilde{\mathrm{D}}(\mathrm{G}_{\alpha})\)，而 \((\tilde{\mathrm{D}}(\mathrm{G}_{\alpha}), \tilde{f}_{\alpha\beta})\) 是满足引理 3 假设的拓扑群射影系统。

由此引理可知，拓扑群 \(\varprojlim\tilde{\mathrm{D}}(\mathrm{G}_{\alpha})\) 同构于一族 \((\mathrm{S}_{\lambda})_{\lambda\in\mathrm{L}}\) 概单紧李群的乘积。由引理 1，同态 \((\pi_{\alpha})\) 所成射影系统的极限可与连续同态 \(\pi:\prod_{\lambda\in\mathrm{L}}\mathrm{S}_{\lambda}\to\overline{\mathrm{D}(\mathrm{G})}\) 识别，该同态是满射（General Topology, Chap. I, §9, no. 6, Cor. 2 of Prop. 8）。

现在注意，群
\[
\prod_ {\lambda \in L} S _ {\lambda}
\]

等于其换位子群：这是由 §4, no. 5, Cor. of Prop. 10 得出的。对于 \(\overline{\mathrm{D(G)}}\)，由于 \(\pi\) 是满射，也同样如此。因此，\(\mathrm{D(G)} \supset \mathrm{D}(\overline{\mathrm{D(G)}}) = \overline{\mathrm{D(G)}}\)。于是群 \(\mathrm{D(G)}\) 是紧的并且等于其换位子群；这证明了 \(a\)，因为关于 \(\mathrm{C(G)}_0\) 的断言是平凡的。

另一方面，\(\pi : \prod_{\lambda \in L} S_{\lambda} \to D(G)\) 的核可与 \(\varprojlim \operatorname{Ker}(\pi_{\alpha})\) 识别（Algebra, Chap. II, §6, no. 1, Remark 1），因而是一个紧的、全不连续的中心子群，故得 \(c\)。

我们证明 \(b\)。对于 I 中的所有 \(\alpha\)，映射 \(s_{\alpha}: \mathrm{C}(\mathrm{G}_{\alpha})_0 \times \mathrm{D}(\mathrm{G}_{\alpha}) \to \mathrm{G}_{\alpha}\) 满足 \(s_{\alpha}(x,y) = xy\)；对于 \(x \in \mathrm{C}(\mathrm{G}_{\alpha})_0, y \in \mathrm{D}(\mathrm{G}_{\alpha})\)，它是满射，且其核是有限中心子群（§1, no. 4, Cor. 1 of Prop. 4）。\(s_{\alpha}\) 构成映射射影系统，其极限根据前述结果可与同态 \((x,y) \mapsto xy\) 从 \(\mathrm{C}(\mathrm{G})_0 \times \mathrm{D}(\mathrm{G})\) 到 \(\mathrm{G}\) 识别。现在如前可知，此映射是满射，且其核是中心的、全不连续的，故得 \(b\)。

推论。每个可解连通紧群都是交换的。

事实上，此时换位子群也是可解的，且等于其换位子群（命题 2 a)），因而退化为单位元。

